\documentclass[12pt]{article}
\usepackage{ceai}


\begin{document}
\title{CE 488/5588 Homework 05\\[2in]
\textbf{Name: \rule{2in}{0.15mm}}}
% \textbf{Name: \FillRule{2in} }
\author{}
\date{}
\maketitle



\newpage
\doublespacing

\section*{Multiple Choice / Fill-in-the-blank Questions}

\subsection*{Question 1} 
\textbf{1a}. Given a logistic regression (binary) classification model, which is expressed as $f[u(\bm{x})]$, the function $u(\bm{x}) = 0$ defines (choose the correct statement below):
\begin{itemize}
    \item [$\square$] the decision boundary function in the space of the input feature vectors.
    \item [$\square$] a linear function (all the time).
    \item [$\square$] a sigmoid function for classification.
    \item [$\square$] the classification model itself.
\end{itemize}

% Your explanation: \FillRule{5in}



\textbf{1b\\}. (Graduate Student) In the above expression, when the discriminative model $f[u(\bm{x})]$ is learned from data, the value of $u(\bm{x})$ measures the probability of a membership (i.e., $P(c|\bm{x})$. Why? 

\FillRule{6.5in}

\FillRule{6.5in}

\FillRule{6.5in}

\FillRule{6.5in}
\newpage



\bigskip

\textbf{Question 2} \\
When learning a logistic regression model $f[u(\bm{x}|\bm{w})]$ from data (each data point $x_k$ has a label 0 or 1, denoted by $c_k$), given that $u(\bm{x}|\bm{w})$ is a linear function, its total loss function is written as:
\[
\mathcal{L}(\bm{w}) = -\sum_k \Big( c_k\ln(f_k) + (1-c_k)\ln(1-f_k) \Big)
\]
What statements below are correct?
\begin{itemize}
    \item [$\square$] $\mathcal{L}$ depicts a surface of $\bm{w}$ that is globally convex with respect to $\bm{w}$.
    \item [$\square$] The total loss function is based on the notion of logarithmic loss.
    \item [$\square$] The gradient of $\mathcal{L}$ with respect to the parameter vector $\bm{w}$ has an analytical closed form.
    \item [$\square$] $c_k = 1$ or $c_k = 0$ acts like a switch, turning off one of the two terms in the loss function.
\end{itemize}

\bigskip

\textbf{Question 3} \\
Consider the binary logistic regression loss for a single data point:
\[
\ell = -c \ln(f) - (1-c)\ln(1-f)
\]
where $f = f[u(x)]$ is the score computed by the binary classifier defined as
\[
c =
\begin{cases}
1, & \text{if } f > \frac{1}{2},\\[1ex]
0, & \text{if } f \leq \frac{1}{2}.
\end{cases}
\]
Perform the following loss calculations (round your answers to two decimal places):
\begin{enumerate}[label=(\arabic*)]
    \item Given $c_k = 1$ and the prob. score $f_k = 0.2$, then:
    \begin{itemize}
           \item The prediction of class label $c^* =$ \FillRule{1in}. 
           \item The loss is \FillRule{1in}.
    \end{itemize}
    \item Given $c_k = 0$ and the prob. score $f_k = 0.8$, then 
        \begin{itemize}
           \item The prediction of class label $c^* =$ \FillRule{1in}. 
           \item The loss is \FillRule{1in}.
    \end{itemize}
\end{enumerate}

\newpage

\section{Programming Practice and Data Operation}
In class, we demonstrated that by using the UCI Machine Learning Concrete dataset (see \href{https://umsystem.instructure.com/courses/280496/files/32742548?module_item_id=9797693}{Canvas Link}), we developed a simple logistic regression model. We took the physical viewpoint that the extracted water-cement ratio might be a good feature to indicate the 28-day concrete strength. The resulting prediction accuracy was low (about $20\%$).

In this work, we will use the original feature values of cement and water contents as bivariate features to make the prediction. This calls for a logistic regression model:
\[
f[u(\bm{x})] = \frac{1}{1+\exp(-w_2 x_2 -w_1 x_1 -w_0)}
\]
i.e., the decision function in the $x_1$--$x_2$ space is a linear function defined by 
\[
-w_2 x_2 -w_1 x_1 -w_0 = 0,
\]
where $x_1$ represents the water feature and $x_2$ represents the cement feature value.

Using the class \href{https://umsystem.instructure.com/courses/280496/files/32742547?module_item_id=9797692}{notebook} as a starting point, resolve the following tasks:
\begin{itemize}
    \item Conduct the logistic regression model as described above, and report the confusion matrix and accuracy based on a training and testing data split.
    \item Report the decision boundary function $-w_2 x_2 -w_1 x_1 -w_0 = 0$ by providing the values of $w_2$, $w_1$, and the intercept $w_0$.
    \item Using a 2D scatter plot, overlay the decision boundary function on the training data. Does your reported accuracy appear consistent with this plot?
\end{itemize}


\end{document}
